% =====================================================================
% Chapter 6 - Conclusion (revised)
%
% Paragraphs 1, 4 and 6 are unchanged from the current draft.
% Paragraph 2 compressed  178 -> 117 words  (dropped the Table 5.1 pointer,
%   the 52-202 range, the 11-of-12 detail and the five-rig-cases count; all
%   are already in Chapter 4 and Table 5.1)
% Paragraph 3 compressed  101 ->  72 words  (the reference-library argument
%   is made at length in Section 5.6; repeating it here was duplication)
% Paragraph 5 prose -> numbered list, 154 -> 164 words, and items 2-4 each
%   gained a sentence saying what the result would establish
%
% NET: -80 counted words against the current draft.
%
% BEFORE PASTING:
%   1. Replace  sec:comparison  with your real Section 5.6 label (2 places).
%   2. \itemsep needs enumitem. If it is not loaded, delete "\itemsep0.3em".
%   3. Keep the existing \chapter{Conclusion} line above this text.
% =====================================================================

One accelerometer on a three-storey shear frame, with screwed joints loosened at four
locations in four grades, supported two analysis routes. Every route that passes through a
parametrised model of the frame fails in the same place. Direct inversion closes every
measured frequency exactly and at two of the three locations where the problem is square,
its solutions require a storey to have stiffened, which loosening a screw cannot do. The one
admissible solution softens a storey the loosened plate does not adjoin. Re-parametrising by
plate does not repair it: no hypothesis fits any measured case and its ranking is right 1 time
in 4, which is chance. The network inherits the storey parametrisation in full. It calls the
bottom or the top storey and almost never the middle one, so fails at Floor~2, whose plate is
the one requiring the middle storey. Its physics term cut the residual it penalises by 17.8\%
and moved no task metric beyond fold and seed scatter on any dataset. Storey stiffness was
recovered at no location. The one route that applies no forward model got every replicated
severe run right. The obstacle lies in the model those routes share rather than in the
identifiability of the inverse problem.

Five teardown and rebuild cycles put the reassembly floors at 0.30, 0.46 and 0.32\%,
measured rather than assumed, and damage at a screwed joint is detectable against them at the
lightest grade tested. Detection is where the two routes agree. The first-mode shift falls
monotonically with grade, but saturates above one turn and the cells below severe are single
application, so severity is an ordinal trend and not a classification. Localisation is the
strongest positive result: every replicated severe run correct, and 18 of 18 at two other
sensor positions against references recorded elsewhere. All of this came from a
consumer-grade part at the noisy end of the published range, with the whole chain running on
the node rather than a base station.

The research question asked whether embedding the equation of motion in a learned inverse
model improves on classical modal analysis. It does not. The classical route is the one to
put on a low-cost node: it returns 6 of 6 where the network returns 3 of 6, and needs no
forward model and no retargeting per structure. Its reference-library requirement, set out in
Section~\ref{sec:comparison}, remains the constraint a working forward model would remove.

The rig is a bolted steel frame and nothing here establishes transfer to unreinforced
masonry, whose failure is governed by strength rather than the stiffness a frequency
observes. No second structure was measured. Five rebuild cycles give an operational floor
rather than a formal exclusion limit, and no thermal variation was imposed, so the floor
bounds reassembly and not deployment. Under these supervised synthetic settings, the physics
term supplied no measurable information beyond the labels. The trained weights behind the rig
predictions were not retained, one of those predictions rests on imputed modes so its
localisation does not stand, and the healthy case is not an independent test. The output is a
change-detection signal. It could prioritise inspection; it is not a safety assessment.

Four experiments follow, in order of what each would settle.

\begin{enumerate}\itemsep0.3em
\item \textbf{A rotational-spring model, fitted to the vectors already recorded.} The
semi-rigid family identified in Section~\ref{sec:comparison} needs no new measurements. A
solution identifies the release the shear chain cannot represent; its absence moves the
discrepancy off the connection.

\item \textbf{Temperature variation at fixed damage.} This replaces the reassembly floor with
an environmental one. The trace-grade margins reported here clear a floor that bounds rebuild
variability only, and would have to be rejudged against it.

\item \textbf{Operational modal analysis under ambient excitation.} No dwelling supplies a
controlled impulse. This establishes whether the identification chain survives without one,
which decides whether the method leaves the laboratory at all.

\item \textbf{The physics-term ablation repeated on unlabelled or heavily noised captures.}
The null reported here rests on labelled synthetic data, where the labels and the inputs came
through the same forward map the residual inverts. This tests the term where it was designed
to operate.
\end{enumerate}

A physics term earns its place where it constrains something the supervision leaves
undetermined, and here the labels and the inputs came through the same forward map the
residual inverts. That bounds when such a term helps, not whether it ever does. In this
controlled campaign, the limit on stiffness recovery came from the structural model rather
than measurement resolution.
